\documentclass[10pt,a4paper]{amsart}
\usepackage[margin=19mm]{geometry}
\usepackage{amsmath,amssymb,mathtools}
\usepackage[scr=rsfs,cal=euler]{mathalfa}
\usepackage{xfrac}
\usepackage[unicode,hidelinks,bookmarks=false]{hyperref}
\newcommand{\ie}{i.e.\ }
\newcommand{\cf}{cf.\ }
\newcommand{\resp}{respectively\ }
\newcommand{\Subsection}{\subsection}
\providecommand{\nobreakdash}{\nobreak\hbox{-}}
\setlength{\emergencystretch}{2em}
\multlinegap0pt
\usepackage{amssymb}
\usepackage{mathtools}
\usepackage[shortlabels]{enumitem}
\newtheorem{theorem}{Theorem}
\newtheorem{corollary}{Corollary}
\DeclareMathOperator{\supp}{supp}
\title{Support-sensitive bounds for shortest zero-sum subsequences}
\author{Claudiu Pop, George C. \c{T}urca\c{s}}
\date{Source: arXiv:2505.04187v2 (CC BY 4.0)}
\begin{document}
\maketitle

\noindent\emph{Source and licence.} Support-sensitive bounds for shortest zero-sum subsequences. Version arXiv:2505.04187v2 (CC BY 4.0), distributed by arXiv under the Creative Commons Attribution 4.0 International licence, http://creativecommons.org/licenses/by/4.0/, as recorded in the arXiv licence metadata for this exact version and shown on the abstract page https://arxiv.org/abs/2505.04187v2. Full licence notice: \texttt{licenses/arxiv-2505.04187-CC-BY-4.0.txt}.\par\smallskip

\noindent\textit{Editorial scope.} Counted unit: the article's corollary cor:support-three-cyclic (source0.tex lines 223--232). The article's complete original proof of it is delivered with it (source0.tex lines 234--284, one \texttt{proof} environment of 51 lines). No mathematics is written, repaired or re-derived: the statement and the proof are the article's own lines, in the article's own notation, and no retained line is substituted or re-flowed.  Retained uncounted as the source-local context the proof needs: lines 139--145; lines 188--193.  Nothing was omitted from any retained span, so the package carries no omission record.  No retained line contains a citation, so the wrapper carries no bibliography and no bibliography entry.  No retained line contains a figure environment, an image-inclusion command or drawing code, so no image is delivered and the package declares no asset dependency.  Editorial formatting only, in addition: an AMS \texttt{amsart} preamble that restates the article's own declarations and macros used by the retained lines, namely the environments \texttt{theorem}, \texttt{corollary} on separate counters, together with the standard packages those lines need.  The retained environments print in this document as Theorem 1, Corollary 1 in order of appearance, whereas the article prints the same items with the numbering of its own full document.  The complete native source of the article as distributed in the arXiv e-print for this exact version is bundled unchanged as \texttt{sources/177726/source0.tex}.
\begin{theorem}\label{thm:finite-abelian-main}
Let $G$ be a finite abelian group of order $n$, and let $S$ be a sequence over $G$ with $|S|=n$.  Then
\[
  MZ(S)\le n-|\supp(S)|+1.
\]
Equivalently, if $MZ(S)=n-s$ for some $s\in\{0,\ldots,n-1\}$, then $|\supp(S)|\le s+1$.
\end{theorem}

\begin{theorem}[Cyclic long-regime refinement]\label{thm:savchev-chen-refinement}
Let $S$ be a length-$n$ sequence over $C_n$.  Put $r=MZ(S)$ and $t=|\supp(S)|$.  If $r-1>n/2$, then
\[
  r\le n-\binom{t}{2}.
\]
\end{theorem}

\begin{corollary}\label{cor:support-three-cyclic}
Let $n\ge 5$, and let $S$ be a length-$n$ sequence over $C_n$ with $|\supp(S)|=3$.  Then
\[
  MZ(S)\le n-3.
\]
This bound is sharp: if $g$ is a generator of $C_n$, then
\[
  MZ\bigl(g^{[n-2]}(2g)(3g)\bigr)=n-3.
\]
\end{corollary}

\begin{proof}
Put $r=MZ(S)$.  By Theorem~\ref{thm:finite-abelian-main},
\[
  r\le n-|\supp(S)|+1=n-2.
\]
Assume first that $n\ge 7$.  If $r>n-3$, then the preceding inequality forces $r=n-2$.  Hence
\[
  r-1=n-3>n/2,
\]
so Theorem~\ref{thm:savchev-chen-refinement}, applied with $|\supp(S)|=3$, gives
\[
  r\le n-\binom{3}{2}=n-3,
\]
a contradiction.  Thus $r\le n-3$ for $n\ge 7$.

It remains to handle $n=5$ and $n=6$.  For $n=5$, if $0\in\supp(S)$, then $MZ(S)=1\le 2=n-3$.  Otherwise $\supp(S)$ is a three-element subset of the four nonzero elements of $C_5$.  These four elements split into two inverse pairs, so any three-element subset contains one inverse pair.  The corresponding two occurrences form a zero-sum subsequence, and $MZ(S)\le 2=n-3$.

For $n=6$, write $C_6=\langle g\rangle$ and identify terms with their coefficients modulo $6$.  If $0\in\supp(S)$, then $MZ(S)=1\le 3=n-3$.  If $\supp(S)$ contains one of the inverse pairs $\{g,5g\}$ or $\{2g,4g\}$, then $S$ has a zero-sum subsequence of length $2$.

We may therefore assume that $\supp(S)$ contains no zero and no inverse pair.  Then it must contain $3g$ and one element from each of the pairs $\{g,5g\}$ and $\{2g,4g\}$.  Thus the support, written in coefficients of $g$, is one of
\[
  \{1,2,3\},\quad \{1,3,4\},\quad \{2,3,5\},\quad \{3,4,5\}.
\]
Let $m_a$ denote the multiplicity of $ag$ in $S$.  If $m_3\ge 2$, then $3g+3g=0$, so $MZ(S)\le 2$.  We may assume $m_3=1$.

For support $\{1,2,3\}$, the subsequence $g(2g)(3g)$ has sum $0$.  For support $\{3,4,5\}$, the subsequence $(3g)(4g)(5g)$ has sum $0$.  For support $\{1,3,4\}$, we have $m_1+m_4=5$.  If $m_1\ge 2$, then $g+g+4g=0$; if $m_1\le 1$, then $m_4\ge 4$, and $4g+4g+4g=0$.  For support $\{2,3,5\}$, we have $m_2+m_5=5$.  If $m_2\ge 3$, then $2g+2g+2g=0$; if $m_2\le 2$, then $m_5\ge 3$, and $2g+5g+5g=0$.  Hence $MZ(S)\le 3=n-3$ in every case.

It remains to prove sharpness.  Let
\[
  S_n=g^{[n-2]}(2g)(3g)
\]
with $g$ a generator of $C_n$.  Since $n\ge 5$, the elements $g,2g,3g$ are distinct.  The subsequence $g^{[n-5]}(2g)(3g)$ has length $n-3$ and sum $ng=0$, so $MZ(S_n)\le n-3$.

Now consider any nonempty subsequence of $S_n$.  It has the form
\[
  g^{[a]}(2g)^{[\epsilon_2]}(3g)^{[\epsilon_3]},
\]
where $0\le a\le n-2$ and $\epsilon_2,\epsilon_3\in\{0,1\}$.  Its length is
\[
  L=a+\epsilon_2+\epsilon_3,
\]
and its coefficient of $g$ is
\[
  c=a+2\epsilon_2+3\epsilon_3=L+\epsilon_2+2\epsilon_3.
\]
If $L<n-3$, then $L\le n-4$.  Since the subsequence is nonempty,
\[
  1\le c\le L+3\le n-1.
\]
Thus $c$ is not divisible by $n$, and the subsequence is not zero-sum.  Therefore no zero-sum subsequence has length less than $n-3$, while one of length $n-3$ exists.  Hence $MZ(S_n)=n-3$.
\end{proof}

\end{document}
